% 原分章PDF111页附录7.I标题—115页附录7.J标题之前，印刷页687—691。
\section*{附录7.I\quad 稀疏迹方法的算法}
\addcontentsline{toc}{section}{附录7.I 稀疏迹方法的算法}
\markright{附录7.I 稀疏迹方法的算法}

首先，利用传统组套索估计矩阵 $A$ 的初始值。实现传统组套索方法的优化问题可定义为
\begin{equation}
\min_{a_{.j}}\quad F=\operatorname{Tr}\bigl((Y-AX)^T(Y-AX)\bigr)
 +\lambda\sum_{j=1}^nm\|a_{.j}\|_2.\tag{7.I.1}
\end{equation}
令
\[
\rho(A)=\operatorname{Tr}\bigl((Y-AX)^T(Y-AX)\bigr).
\]
利用矩阵微积分，得到
\begin{equation}
\frac{\partial\rho(A)}{\partial A}=-(Y-AX)X^T,\tag{7.I.2}
\end{equation}
进一步化为
\begin{equation}
\frac{\partial\rho(A)}{\partial A}
 =-(Y-A^{(-j)}X)X^T+a_{.j}\|X_{j.}\|_2^2,\tag{7.I.3}
\end{equation}
其中
\[
A^{(-j)},X_{j.}
\]
定义于结论7.22。由式（7.I.3），有
\begin{equation}
\frac{\partial\rho(A)}{\partial a_{.j}}
 =-(Y-A^{(-j)}X)X_{j.}^T+a_{.j}\|X_{j.}\|_2^2.\tag{7.I.4}
\end{equation}
众所周知，范数 $\|a_{.j}\|_2$ 的次梯度为
\begin{equation}
\partial\|a_{.j}\|_2=\begin{cases}
\dfrac{a_{.j}}{\|a_{.j}\|_2}&\|a_{.j}\|_2\ne0,\\
s,\ \|s\|_2\le1&\|a_{.j}\|_2=0.
\end{cases}\tag{7.I.5}
\end{equation}
点 $a_{.j}^*$ 为式（7.I.1）中目标函数 $F$ 的极小点，当且仅当
\begin{equation}
0\in\partial F(a_{.j}^*)=\frac{\partial\rho(A)}{\partial a_{.j}}
 +\lambda m\,\partial\|a_{.j}\|_2.\tag{7.I.6}
\end{equation}
由式（7.I.5）、（7.I.6），当 $\|a_{.j}\|_2\ne0$ 时，有
\begin{equation}
-(Y-A^{(-j)}X)X_{j.}^T+a_{.j}\|X_{j.}\|_2^2
 +\frac{\lambda ma_{.j}}{\|a_{.j}\|_2}=0.\tag{7.I.7}
\end{equation}
令
\begin{equation}
S_j=(Y-A^{(-j)}X)X_{j.}^T.\tag{7.I.8}
\end{equation}
将式（7.I.8）代入式（7.I.7），得到
\begin{equation}
-S_j+a_{.j}\|X_j\|_2^2+\frac{\lambda ma_{.j}}{\|a_{.j}\|_2}=0,\tag{7.I.9}
\end{equation}
从而
\begin{equation}
\left(\|X_{j.}\|_2^2+\frac{\lambda m}{\|a_{.j}\|_2}\right)a_{.j}=S_j.\tag{7.I.10}
\end{equation}
对式（7.I.10）两边取范数 $\|\cdot\|_2$，得到
\[
\left(\|X_{j.}\|_2^2+\frac{\lambda m}{\|a_{.j}\|_2}\right)
 \|a_{.j}\|_2=\|S_j\|_2,
\]
即
\begin{equation}
\|X_{j.}\|_2^2\|a_{j.}\|_2+\lambda m=\|S_j\|_2.\tag{7.I.11}
\end{equation}
由式（7.I.11）解出 $\|a_{j.}\|_2$，得到
\begin{equation}
\|a_{j.}\|_2=\frac{\|S_j\|_2-\lambda m}{\|X_{j.}\|_2^2}.\tag{7.I.12}
\end{equation}
将式（7.I.12）代入式（7.I.10），得到
\begin{equation}
a_{.j}=\left(1-\frac{\lambda m}{\|S_j\|_2}\right)
 \frac{S_j}{\|X_{j.}\|_2^2}.\tag{7.I.13}
\end{equation}
现在考虑 $\|a_{.j}\|_2=0$ 的情形。此时式（7.I9）变为
\begin{equation}
-S_j+\lambda m\delta=0,\tag{7.I.14}
\end{equation}
其中 $\|\delta\|_2\le1$，即
\[
\|S_j\|_2=\lambda m\|\delta\|_2\le\lambda m,
\]
从而
\begin{equation}
\frac{\lambda m}{\|S_j\|_2}\ge1.\tag{7.I.15}
\end{equation}
换言之，若 $\lambda m/\|S_j\|_2\ge1$，则必有
\begin{equation}
a_{.j}=0,\qquad\text{当}\quad\frac{\lambda m}{\|S_j\|_2}\ge1.\tag{7.I.16}
\end{equation}
结合式（7.I.13）、（7.I.16），得到最终解：
\begin{equation}
a_{.j}=\left(1-\frac{\lambda m}{\|S_j\|_2}\right)_+
 \frac{S_j}{\|X_{j.}\|_2^2},\tag{7.I.17}
\end{equation}
其中
\[
(b)_+=\begin{cases}b&b>0,\\0&b\le0.\end{cases}
\]
接下来研究算法中的主要步骤。类似于式（7.I.4），有
\begin{equation}
\frac{\partial\rho(A)}{\partial a_{kj}}
 =-(Y_k-a_k^{-(kj)}X)X_j^T+a_{kj}\|X_j\|_2^2,\tag{7.I.18}
\end{equation}
其中
\[
Y_{k.}=\begin{bmatrix}y_{k1}&\cdots&y_{kL}\end{bmatrix},\qquad
a_{k.}=\lfloor a_{k1}\quad\cdots\quad a_{kj}\quad\cdots a_{kn}\rfloor,
\]
而 $a_{k.}^{(kj)}$ 是将 $a_{k.}$ 中的第 $kj$ 个元素替换为零所得的向量。令 $S_{kj}=(Y_{k.}-a_{k.}^{-(kj)}X)X_{j.}^T$，式（7.I.18）可改写为
\begin{equation}
\frac{\partial\rho(A)}{\partial a_{kj}}=-S_{kj}+a_{kj}\|X_{j.}\|_2^2.\tag{7.I.19}
\end{equation}
同样，$\|a_{.j}\|_2$ 在点 $a_{kj}$ 处的次微分为
\begin{equation}
\frac{\partial\|a_{.j}\|_2}{\partial a_{kj}}
 =\begin{cases}
\dfrac{a_{kj}}{\|a_{.j}\|_2}&a_{kj}\ne0,\\
|\alpha|\le1&a_{kj}=0.
\end{cases}\tag{7.I.20}
\end{equation}
现在考虑 $a_{kj}\ne0$。此时，求解问题（7.144）的最优条件为
\begin{equation}
\frac{\partial F}{\partial a_{kj}}
 =-S_{kj}+a_{kj}\|X_{j.}\|_2^2
 +\lambda\left(\frac n{\|a_{k.}\|_2}+\frac m{\|a_{.j}\|_2}\right)a_{kj}=0.\tag{7.I.21}
\end{equation}
由式（7.I.21）解出 $a_{kj}$，得到
\begin{equation}
a_{kj}^{(t+1)}=\frac{S_{kj}^{(t)}}
 {\|X_{j.}\|_2^2+\lambda\left(\dfrac n{\|a_{k.}^{(t)}\|_2}
 +\dfrac m{\|a_{.j}^{(t)}\|_2}\right)},\tag{7.I.22}
\end{equation}
其中 $S_{kj}^{(t)},a_{k.}^{(t)},a_{.j}^{(t)}$ 均由矩阵 $A^{(t)}$ 计算。接下来考虑 $a_{kj}=0$。此时式（7.I.21）变为
\begin{equation}
\frac{\partial F}{\partial a_{kj}}=-S_{kj}+\lambda(m+n)\alpha=0,\tag{7.I.23}
\end{equation}
其中
\[
|\alpha|\le1.
\]
因此，若 $|S_{kj}|\le\lambda(m+n)$，则令 $a_{kj}=0$。
